% Generado por scripts/thesis_artifacts.py - no editar a mano.
\chapter{Resultados}

\section{Resumen}

Se compararon las dos microarquitecturas RV32IM sobre la FPGA Cyclone V DE1-SoC bajo el protocolo descrito en el capítulo de metodología. El pipeline alcanzó una Fmax de 54.45 MHz ($\pm$ 0.78 MHz, $n = 5$) frente a los 52.92 MHz ($\pm$ 0.73 MHz, $n = 5$) del monociclo, una mejora de frecuencia del 2.9 \%.

\section{Verificación funcional}

La corrección funcional de ambas implementaciones se verificó contra el modelo de referencia de Python y contra la suite \texttt{riscv-tests} antes de medir el rendimiento. Los resultados se resumen en la tabla \ref{tab:verification_results}.

% Generado por scripts/thesis_artifacts.py - no editar a mano.
% Fecha: 2026-10-21 16:38
\begin{table}[htbp]
\centering
\caption{Resultados de la verificación funcional (suites de instrucciones, modelo de referencia y hazards).}
\label{tab:verification_results}
\begin{tabular}{lrrrrrr}
\toprule
Suite de pruebas & \multicolumn{3}{c}{Monociclo} & \multicolumn{3}{c}{Pipeline} \\
\cmidrule(lr){2-4} \cmidrule(lr){5-7}
& Aprobadas & Fallidas & Total & Aprobadas & Fallidas & Total \\
\midrule
RV32I ISA (37 tests) & 37 & 0 & 37 & \textemdash & \textemdash & \textemdash \\
RV32M M-Extension (8 tests) & 8 & 0 & 8 & \textemdash & \textemdash & \textemdash \\
RV32MI Machine-Mode (15 tests) & 15 & 0 & 15 & \textemdash & \textemdash & \textemdash \\
ALU Unit Tests & 27 & 0 & 27 & \textemdash & \textemdash & \textemdash \\
Branch Unit Tests & 20 & 0 & 20 & \textemdash & \textemdash & \textemdash \\
Model vs DUT Comparison & 4 & 0 & 4 & \textemdash & \textemdash & \textemdash \\
test\_validation & 1 & 0 & 1 & 1 & 0 & 1 \\
\bottomrule
\end{tabular}
\end{table}


\section{Frecuencia máxima y recursos}

% Generado por scripts/thesis_artifacts.py - no editar a mano.
% Fecha: 2026-10-21 16:16
\begin{table}[htbp]
\centering
\caption{Frecuencia máxima por réplica de síntesis (modelo \emph{Slow 1100 mV 85$^\circ$C}).}
\label{tab:fmax_replicas}
\begin{tabular}{lrr}
\toprule
Réplica & Monociclo (MHz) & Pipeline (MHz) \\
\midrule
1 & 53.75 & 54.59 \\
2 & 52.89 & 54.19 \\
3 & 51.75 & 54.51 \\
4 & 53.19 & 55.57 \\
5 & 53.02 & 53.41 \\
\midrule
Media & 52.92 & 54.45 \\
Desviación estándar & 0.73 & 0.78 \\
Mínimo & 51.75 & 53.41 \\
Máximo & 53.75 & 55.57 \\
Número de réplicas (n) & 5 & 5 \\
\bottomrule
\end{tabular}
\end{table}


% Generado por scripts/thesis_artifacts.py - no editar a mano.
% Fecha: 2026-10-21 10:55
\begin{table}[htbp]
\centering
\caption{Utilización de recursos de la FPGA por réplica (media $\pm$ desviación estándar, n = 5).}
\label{tab:resources}
\begin{tabular}{lrrl}
\toprule
Recurso & Monociclo & Pipeline & Overhead \\
\midrule
ALMs & \textemdash & \textemdash & \textemdash \\
Registros & \textemdash & \textemdash & \textemdash \\
Bloques RAM & \textemdash & \textemdash & \textemdash \\
Bloques DSP & \textemdash & \textemdash & \textemdash \\
Bits de memoria & \textemdash & \textemdash & \textemdash \\
Pines & \textemdash & \textemdash & \textemdash \\
\bottomrule
\end{tabular}
\end{table}


% Generado por scripts/thesis_artifacts.py - no editar a mano.
\begin{figure}[htbp]
\centering
\begin{tikzpicture}
\begin{axis}[
  width=0.65\textwidth,
  height=7cm,
  ybar,
  bar width=25pt,
  ylabel={Fmax (MHz)},
  xtick={1,2},
  xticklabels={Monociclo,Pipeline},
  nodes near coords,
  nodes near coords align={vertical},
  ymin=0,
  error bars/.cd, y dir=both, y explicit,
]
  \addplot coordinates {(1,52.92) +- (0,0.73), (2,54.45) +- (0,0.78)};
\end{axis}
\end{tikzpicture}
\caption{Frecuencia máxima de cada microarquitectura, media $\pm$ desviación estándar de las cinco réplicas de síntesis.}
\label{fig:fmax}
\end{figure}


\section{CPI, IPC y throughput}

% Generado por scripts/thesis_artifacts.py - no editar a mano.
% Fecha: 2026-10-21 10:55
\begin{table}[htbp]
\centering
\small
\caption{CPI, IPC, throughput y latencia por programa y arquitectura. El throughput y la latencia se obtienen con la Fmax media de las cinco réplicas de síntesis.}
\label{tab:benchmarks}
\begin{tabular}{llrrrrrr}
\toprule
Programa & Arquitectura & Ciclos & Instrucciones & CPI & IPC & Throughput (MIPS) & Latencia (ms) \\
\midrule
branch\_stress & Monociclo & 4004 & 4004 & 1.0000 & 1.0000 & 36.91 & 0.1085 \\
branch\_stress & Pipeline & 8005 & 4003 & 1.9998 & 0.5001 & 28.31 & 0.1414 \\
coremark\_iter1 & Monociclo & 327888 & 327854 & 1.0001 & 0.9999 & 36.91 & 8.8825 \\
coremark\_iter1 & Pipeline & 457786 & 327372 & 1.3984 & 0.7151 & 40.48 & 8.0861 \\
coremark\_iter10 & Monociclo & 3152076 & 3152042 & 1.0000 & 1.0000 & 36.91 & 85.3897 \\
coremark\_iter10 & Pipeline & 4417297 & 3151560 & 1.4016 & 0.7135 & 40.39 & 78.0248 \\
load\_use\_stress & Monociclo & 1125 & 1125 & 1.0000 & 1.0000 & 36.91 & 0.0305 \\
load\_use\_stress & Pipeline & 2236 & 1124 & 1.9893 & 0.5027 & 28.46 & 0.0395 \\
m\_ext\_stress & Monociclo & 1175 & 155 & 7.5806 & 0.1319 & 4.87 & 0.0318 \\
m\_ext\_stress & Pipeline & 1236 & 154 & 8.0260 & 0.1246 & 7.05 & 0.0218 \\
raw\_stress & Monociclo & 1148 & 1148 & 1.0000 & 1.0000 & 36.91 & 0.0311 \\
raw\_stress & Pipeline & 1435 & 1147 & 1.2511 & 0.7993 & 45.25 & 0.0253 \\
\bottomrule
\end{tabular}
\end{table}


El monociclo presenta CPI = 1 en los programas sin división: cada ciclo retira una instrucción. La excepción es el programa de la extensión M, donde el divisor de 34 ciclos mantiene el datapath ocupado y el CPI del monociclo sube a 7.58. El pipeline retira una instrucción por ciclo en el mejor caso y pierde ciclos por hazard de datos, por vaciado de rama y por la latencia de la división.

% Generado por scripts/thesis_artifacts.py - no editar a mano.
\begin{figure}[htbp]
\centering
\begin{tikzpicture}
\begin{axis}[
  width=0.95\textwidth,
  height=7cm,
  ybar,
  bar width=8pt,
  ylabel={Throughput (MIPS)},
  xtick={1,2,3,4,5,6},
  xticklabels={{branch\_stress}, {coremark\_iter1}, {coremark\_iter10}, {load\_use\_stress}, {m\_ext\_stress}, {raw\_stress}},
  xmin=0.5,
  enlarge x limits=0.1,
  legend style={at={(0.5,-0.25)},anchor=north,legend columns=2},
]
  \addplot coordinates {(1,52.920), (2,52.915), (3,52.920), (4,52.920), (5,6.981), (6,52.920)};
  \addplot coordinates {(1,27.230), (2,38.940), (3,38.851), (4,27.373), (5,6.785), (6,43.525)};
  \legend{Monociclo,Pipeline}
\end{axis}
\end{tikzpicture}
\caption{Throughput efectivo por programa y microarquitectura ($F_{\max}/\mathrm{CPI}$).}
\label{fig:throughput}
\end{figure}


% Generado por scripts/thesis_artifacts.py - no editar a mano.
\begin{figure}[htbp]
\centering
\begin{tikzpicture}
\begin{axis}[
  width=0.95\textwidth,
  height=7cm,
  ybar,
  bar width=8pt,
  ylabel={CPI},
  xtick={1,2,3,4,5,6},
  xticklabels={{branch\_stress}, {coremark\_iter1}, {coremark\_iter10}, {load\_use\_stress}, {m\_ext\_stress}, {raw\_stress}},
  xmin=0.5,
  enlarge x limits=0.1,
  legend style={at={(0.5,-0.25)},anchor=north,legend columns=2},
]
  \addplot coordinates {(1,1.000), (2,1.000), (3,1.000), (4,1.000), (5,7.581), (6,1.000)};
  \addplot coordinates {(1,2.000), (2,1.398), (3,1.402), (4,1.989), (5,8.026), (6,1.251)};
  \legend{Monociclo,Pipeline}
\end{axis}
\end{tikzpicture}
\caption{CPI efectivo por programa y microarquitectura.}
\label{fig:cpi}
\end{figure}


\section{Análisis estadístico}

Con $n = 5$ réplicas por tratamiento, el análisis es descriptivo. La desviación estándar de la Fmax es de 0.73 MHz en el monociclo (CV 1.38 \%) y de 0.78 MHz en el pipeline (CV 1.43 \%).

% Generado por scripts/thesis_artifacts.py - no editar a mano.
% Fecha: 2026-10-21 16:38
\begin{table}[htbp]
\centering
\caption{Estadística descriptiva de las cinco réplicas de síntesis por arquitectura (desviación estándar muestral, $n-1$).}
\label{tab:statistics}
\begin{tabular}{llrrrr}
\toprule
Arquitectura & Métrica & Media & Desv. estándar & CV (\%) & Mínimo & Máximo \\
\midrule
Monociclo & Fmax (MHz) & 52.920 & 0.732 & 1.38 & 51.750 & 53.750 \\
Monociclo & ALMs & 11400.000 & 34.066 & 0.30 & 11362.000 & 11438.000 \\
Monociclo & Registros & 18896.800 & 16.084 & 0.09 & 18880.000 & 18917.000 \\
Monociclo & Bloques RAM & 31.000 & 0.000 & 0.00 & 31.000 & 31.000 \\
Monociclo & Bloques DSP & 7.000 & 0.000 & 0.00 & 7.000 & 7.000 \\
Monociclo & Bits de memoria & 229888.000 & 0.000 & 0.00 & 229888.000 & 229888.000 \\
Monociclo & Pines & 85.000 & 0.000 & 0.00 & 85.000 & 85.000 \\
Pipeline & Fmax (MHz) & 54.454 & 0.779 & 1.43 & 53.410 & 55.570 \\
Pipeline & ALMs & 3633.600 & 8.620 & 0.24 & 3623.000 & 3647.000 \\
Pipeline & Registros & 3909.600 & 11.567 & 0.30 & 3895.000 & 3921.000 \\
Pipeline & Bloques RAM & 35.000 & 0.000 & 0.00 & 35.000 & 35.000 \\
Pipeline & Bloques DSP & 10.000 & 0.000 & 0.00 & 10.000 & 10.000 \\
Pipeline & Bits de memoria & 246272.000 & 0.000 & 0.00 & 246272.000 & 246272.000 \\
Pipeline & Pines & 85.000 & 0.000 & 0.00 & 85.000 & 85.000 \\
\bottomrule
\end{tabular}
\end{table}


\section{Eficiencia relativa}

% Generado por scripts/thesis_artifacts.py - no editar a mano.
% Fecha: 2026-10-21 10:55
\begin{table}[htbp]
\centering
\caption{Eficiencia relativa: throughput por recurso utilizado y overhead de recursos del pipeline.}
\label{tab:efficiency}
\begin{tabular}{lrr}
\toprule
Indicador & Monociclo & Pipeline \\
\midrule
Throughput por 1000 ALMs (MIPS) & --- & --- \\
Throughput por 1000 registros (MIPS) & --- & --- \\
\bottomrule
\end{tabular}
\end{table}


% Generado por scripts/thesis_artifacts.py - no editar a mano.
\begin{figure}[htbp]
\centering
\begin{tikzpicture}
\begin{axis}[
  width=0.95\textwidth,
  height=7cm,
  ybar,
  bar width=8pt,
  ylabel={Cantidad de recursos},
  xtick={1,2,3,4},
  xticklabels={{ALMs}, {Registros}, {RAM}, {DSP}},
  xmin=0.5,
  enlarge x limits=0.1,
  legend style={at={(0.5,-0.25)},anchor=north,legend columns=2},
]
  \addplot coordinates {(1,11400.000), (2,18896.800), (3,31.000), (4,7.000)};
  \addplot coordinates {(1,3633.600), (2,3909.600), (3,35.000), (4,10.000)};
  \legend{Monociclo,Pipeline}
\end{axis}
\end{tikzpicture}
\caption{Utilización de recursos de la FPGA por microarquitectura.}
\label{fig:resources}
\end{figure}


\section{Evaluación de la hipótesis}

% Generado por scripts/thesis_artifacts.py - no editar a mano.
% Fecha: 2026-10-21 10:55
\begin{table}[htbp]
\centering
\caption{Speedup del pipeline respecto al monociclo y evaluación de la hipótesis por programa.}
\label{tab:speedup}
\begin{tabular}{lrrrl}
\toprule
Programa & Speedup medio & Desviación estándar & Media $-$ SD & Resultado \\
\midrule
branch\_stress & 0.767 & 0.030 & 0.737 & H0 no descartada \\
coremark\_iter1 & 1.097 & 0.044 & 1.053 & H1 sostenida \\
coremark\_iter10 & 1.094 & 0.043 & 1.051 & H1 sostenida \\
load\_use\_stress & 0.771 & 0.031 & 0.740 & H0 no descartada \\
m\_ext\_stress & 1.449 & 0.058 & 1.391 & H1 sostenida \\
raw\_stress & 1.226 & 0.049 & 1.177 & H1 sostenida \\
\bottomrule
\end{tabular}
\end{table}


\section{Limitaciones}

\begin{itemize}
\item El tamaño muestral ($n = 5$) limita el análisis a la estadística descriptiva.
\item El CPI se mide en simulación; la única fuente de dispersión considerada es la variabilidad de la Fmax entre réplicas.
\item CoreMark se mide con diez iteraciones; la mezcla de instrucciones, y por tanto el CPI, es la de la corrida completa.
\item CoreMark no cabe en la memoria de instrucciones de la FPGA, por lo que no se midió en hardware.
\end{itemize}
